\chapter{Spreads sur matières premières}\label{ch:s2:commodity-spreads}

Une raffinerie gagne l'écart entre le brut qu'elle achète et les produits qu'elle vend. Un trader qui détient cet écart, plutôt que du pétrole, parie sur
les marges de raffinage, le stockage et les saisons, et non sur le cours du pétrole. De 2006 à 2026, le crack spread 3-2-1 de New York Harbor, calculé à partir des cours spot
de l'EIA, s'est établi en moyenne à 21.30~\$ le baril et est revenu à la moyenne avec une demi-vie d'environ 51 jours de bourse. Une règle qui parie sur le retour à la moyenne du spread a obtenu un ratio de Sharpe de 0.48
avec une corrélation de $-0.02$ avec les variations quotidiennes du brut. Sur le marché synthétique de ce chapitre, où le crack présente une saisonnalité programmée et un retour
à la moyenne programmé, la même règle obtient un ratio de Sharpe de 1.05, avec une corrélation de $-0.12$ avec le brut. La construction porte sur \texttt{firm.commspread}.

\section{Calendar spreads, à nouveau}

Un calendar spread (livre~1, chapitre~19) est le plus simple des spreads sur matières premières~: une même matière première à deux échéances, évaluée par le coût de stockage et le rendement d'opportunité (convenience yield). Les
spreads de ce chapitre relient des matières premières, des qualités ou des lieux différents. Chacun est évalué par une activité physique qui transforme l'une en l'autre~: le raffinage,
la production d'électricité, le broyage, le transport maritime, le mélange. Cette activité borne le spread. Quand le crack est trop étroit, les raffineurs réduisent leurs cadences de traitement et les stocks de produits
baissent~; quand il est trop large, ils accélèrent. C'est la raison économique du retour à la moyenne que parie le trader, et la raison pour laquelle ce pari peut échouer quand
l'activité elle-même est bloquée.

\begin{definition}[Trade de marge de transformation]\label{def:s2:commodity-spreads:processing}
Un \emph{trade de marge de transformation}\index{trade de marge de transformation} (processing-spread trade) est une position sur une matière première d'entrée contre les produits qui en sont tirés, dans le ratio que donne le
procédé (brut contre essence et gazole, gaz contre électricité à un heat rate, soja contre tourteau et huile), qui profite des variations de la
marge de transformation plutôt que du niveau des cours des matières premières.
\end{definition}

Girma et Paulson ont constaté que les cours des contrats à terme sur pétrole brut, sur essence et sur fioul domestique sont cointégrés, les spreads entre le brut et ses produits étant stationnaires. Ils
ont relevé, sur le passé, des arbitrages de risque rentables et statistiquement significatifs sur trois spreads pétroliers courants, sans pouvoir être certains que
ces opportunités existent encore aujourd'hui. \texttt{firm.commspread} construit les spreads par ratio
(le \cref{lst:s2:commodity-spreads:build})~; avec une essence à 2.50~\$ le gallon, un gazole à 3.00~\$ et un brut à 80~\$ le baril, le crack 3-2-1 vaut 32~\$.

\section{Qualité et localisation}

\begin{definition}[Spread géographique]\label{def:s2:commodity-spreads:location}
Un \emph{spread géographique}\index{spread géographique} (location spread) est l'écart de prix entre une même matière première, ou une matière presque identique, en deux points de livraison, borné
en temps normal par le coût de son transport entre ces points et élargi sans limite lorsque la capacité de transport est saturée.
\end{definition}

\begin{definition}[Spread de qualité]\label{def:s2:commodity-spreads:quality}
Un \emph{spread de qualité}\index{spread de qualité} (quality spread) est l'écart de prix entre deux qualités d'une même matière première (brut léger et doux contre brut lourd et soufré,
blé à haute teneur en protéines contre blé à basse teneur), fixé par le coût de transformation de la qualité inférieure en produits de la qualité supérieure.
\end{definition}

Le Brent est un brut transporté par mer, dont le prix est fixé en mer du Nord~; le WTI est livré à l'intérieur des terres, à Cushing, dans l'Oklahoma. De 2006 à 2026, le Brent s'est négocié en moyenne 4.76~\$
au-dessus du WTI, avec un écart type de 5.97~\$ (la \cref{fig:s2:commodity-spreads:real}). L'EIA a attribué le spread record de 2011, quand le Brent
valait en moyenne environ 111~\$ et le WTI environ 95~\$, à des goulets d'étranglement dans le transport près de Cushing. En 2017, l'EIA écrivait que les variations à court terme du spread
proviennent généralement de changements de capacité des pipelines ou de la production américaine de brut. Un \omterm{def:s2:commodity-spreads:location}{spread géographique} est un pari sur les pipelines.

Les données quotidiennes contiennent deux mises en garde. Le 20 avril 2020, le cours spot du WTI était de $-\$36.98$, et le Brent se situait 54.34~\$ au-dessus, soit le maximum de la série~: un point
de livraison qui ne pouvait plus accepter de brut l'a évalué sous zéro. Par ailleurs, les variations quotidiennes du spread Brent-WTI présentent une autocorrélation de $-0.36$. Le Brent est
évalué à la clôture de Londres et le WTI à celle de New York, de sorte que le spread de chaque jour mélange deux heures différentes, et une partie de sa variation est annulée le lendemain. Une règle qui
parie sur le retour des écarts du spread vers sa moyenne glissante a obtenu un ratio de Sharpe de 0.85 hors avril 2020. Entrée un jour plus tard, elle n'obtient que 0.32.

\begin{omfigure}
\begin{tikzpicture}
\begin{axis}[width=11cm,height=5.4cm,xlabel={\small année},ylabel={\small \$/baril, moyenne mensuelle},tick label style={font=\footnotesize},
  xmin=2006,xmax=2027,ymin=-10,ymax=70,grid=major,grid style={gray!25},table/col sep=comma,
  x tick label style={/pgf/number format/1000 sep={}},legend style={font=\scriptsize,at={(0.5,0.97)},anchor=north}]
\addplot[omThm,thick] table[x=year,y=crack]{figdata/strategies-2/21-commodity-spreads/real.csv};
\addplot[omDef,thick] table[x=year,y=brentwti]{figdata/strategies-2/21-commodity-spreads/real.csv};
\legend{crack 3-2-1 de New York Harbor,Brent moins WTI}
\end{axis}
\end{tikzpicture}

\omcaption{Moyennes mensuelles du crack spread 3-2-1 de New York Harbor et du spread Brent moins WTI, à partir des cours spot quotidiens de l'EIA, de juin 2006 à septembre 2026.
Données~: EIA via FRED~; \texttt{s2\_fetch\_spreads.py}.}\label{fig:s2:commodity-spreads:real}
\end{omfigure}

\section{Spreads de transformation~: crack, spark, crush}

Le crack synthétique de \texttt{firm.commspread} comporte quatre éléments~: une moyenne de 20~\$, un cycle saisonnier d'amplitude 5~\$ qui culmine fin avril, avant la
saison des grands déplacements estivaux, une déviation qui revient à la moyenne avec une demi-vie de 40 jours, et un retard programmé~: quand le brut bondit, les produits suivent sur plusieurs semaines, de sorte que le
crack se resserre d'abord. Ses variations quotidiennes présentent une corrélation de $-0.36$ avec celles du brut, proche du $-0.38$ du crack 3-2-1 réel. Pari sur le retour du crack
au-delà de 1.5 écart type de sa moyenne glissante sur un an (le \cref{lst:s2:commodity-spreads:fade})~:

\begin{center}
\small
\begin{tabular}{@{}lccc@{}}
\toprule
neuf ans, 0.10~\$/baril par unité échangée & \$/baril par an & ratio de Sharpe & corrélation avec le brut\\
\midrule
parier sur le retour du crack & 15.4 & 1.05 & $-0.12$\\
parier sur le retour du crack, saison retirée & 6.2 & 0.44 & $-0.07$\\
acheter le crack de février à avril & 7.4 & 0.78 & $-0.15$\\
parier sur le retour du \omterm{def:s2:commodity-spreads:location}{spread géographique} & 3.8 & 0.48 & 0.01\\
parier sur le retour du \omterm{def:s2:commodity-spreads:quality}{spread de qualité} & 2.6 & 0.98 & 0.01\\
combinaison à risque égal & & 1.69 & $-0.13$\\
\bottomrule
\end{tabular}
\end{center}

Retirer la moyenne de chaque mois avant de parier sur le retour réduit de plus de moitié le gain de la stratégie sur le crack~: une grande part de ce que la règle simple exploite est la saison. Les deux
règles sur le crack ne forment qu'un seul pari sur le même cycle saisonnier. Sur le crack de l'essence réel, la variation du début de février à la fin d'avril a été positive dans 14
des 20 années de 2007 à 2026, avec une moyenne de 6.17~\$ et une statistique $t$ de 2.8~; la pire année a perdu 11.97~\$.

Le spark spread est l'électricité moins le gaz au heat rate~: à 50~\$ le mégawattheure, un gaz à 3~\$ le million de British thermal units et un heat rate de 7, il
vaut 29~\$. Un trader qui anticipe quelles centrales fixeront le prix, et à quel heat rate, négocie le spread à ce heat rate. Le crush est le tourteau de soja
et l'huile moins les graines dont ils sont issus. Simon a constaté que les écarts du crush par rapport à son équilibre de long terme étaient transitoires de 1985 à 1995, avec
une saisonnalité marquée et une tendance à revenir vers sa moyenne des cinq derniers jours, et que des règles de trading fondées là-dessus auraient été rentables.

\section{Risque dans les livres de spreads}

Un spread est moins volatil que ses jambes, et un livre de spreads semble bien diversifié~: la combinaison ci-dessus affiche un ratio de Sharpe de 1.69. Trois risques
s'y cachent. D'abord, le retour à la moyenne des spreads dépend de l'activité physique qui relie les jambes, et lorsque cette activité est bloquée, le spread part dans une tendance.
Le \omterm{def:s2:commodity-spreads:location}{spread géographique} synthétique a rapporté 35.3~\$ le baril sur quatre ans en pariant sur le retour à la moyenne, puis un goulet d'étranglement programmé sur les pipelines l'a élargi de 15~\$ en six mois.
La stratégie a perdu 14.1~\$ pendant ces mois et 9.3~\$ sur toute la durée du goulet, avant de gagner 8.4~\$ ensuite
(la \cref{fig:s2:commodity-spreads:books}).

Deuxièmement, un spread représente deux positions ou plus sur contrats à terme, chacune faisant l'objet d'un appel de marge, chacune ayant sa propre liquidité, sa livraison et son échéance~; c'est la jambe la moins liquide qui fixe le coût.
Troisièmement, le ratio retenu dans la transaction est une convention, et non le procédé~: les rendements d'une raffinerie diffèrent du 3-2-1, et les 44 et 11 livres du crush sont des moyennes.
Un livre qui parie sur le retour de plusieurs spreads partage aussi leurs facteurs communs. Sur les données spot de l'EIA, la stratégie sur le crack et celle sur le spread Brent-WTI ont toutes deux eu leur meilleure journée
de l'échantillon les 20 ou 21 avril 2020, quand le cours du WTI est devenu négatif~; durant une semaine de ce type, un livre de spreads ne forme qu'une seule position.

\begin{omfigure}
\begin{tikzpicture}
\begin{axis}[width=11cm,height=5.4cm,xlabel={\small année},ylabel={\small cumulé \$/baril par unité},tick label style={font=\footnotesize},
  xmin=1,xmax=10,ymin=-10,ymax=150,grid=major,grid style={gray!25},table/col sep=comma,
  legend style={font=\scriptsize,at={(0.5,-0.3)},anchor=north},legend columns=2]
\fill[gray!15] (axis cs:5,-10) rectangle (axis cs:6.17,150);
\addplot[omThm,thick] table[x=year,y=crack]{figdata/strategies-2/21-commodity-spreads/books.csv};
\addplot[gray,thick,dashed] table[x=year,y=seasonal]{figdata/strategies-2/21-commodity-spreads/books.csv};
\addplot[omDef,thick] table[x=year,y=location]{figdata/strategies-2/21-commodity-spreads/books.csv};
\addplot[black,thick] table[x=year,y=quality]{figdata/strategies-2/21-commodity-spreads/books.csv};
\legend{parier sur le retour du crack,crack de février--avril,parier sur le retour du spread géographique,parier sur le retour du spread de qualité}
\end{axis}
\end{tikzpicture}

\omcaption{P\&L cumulé de quatre règles de spreads synthétiques, en dollars par baril et par unité de spread~; la bande grisée est le goulet d'étranglement du pipeline
introduit qui élargit le \omterm{def:s2:commodity-spreads:location}{spread géographique}. Données~: \texttt{s2\_commspread.market}.}\label{fig:s2:commodity-spreads:books}
\end{omfigure}

\section{Fiches de stratégie}

\begin{strategyfile}{Retour à la moyenne du crack}\label{strat:s2:commodity-spreads:crack}
\sfield{Qui paie, et pourquoi} Les raffineurs et les acheteurs de produits qui couvrent leurs positions aux extrêmes~; les variations de run des raffineurs referment l'écart avec un décalage.
\sfield{Instruments et plateformes} Contrats à terme sur brut, essence et gazole dans le ratio 3-2-1~; contrats de crack spread cotés en bourse.
\sfield{Signal} L'écart du crack par rapport à sa moyenne glissante, avec ou sans sa saisonnalité.
\sfield{Dimensionnement et exécution} Négocié comme un seul ordre de spread lorsque la bourse le permet~; dimensionné par la volatilité du spread.
\sfield{Coûts} Trois jambes~; la plus étroite fixe le coût.
\sfield{Comment elle meurt} Les pannes de raffinerie et les ruptures structurelles de la demande de produits qui déplacent la moyenne.
\sfield{Horizon, capacité, infrastructure} De quelques semaines à quelques mois.
\sfield{Backtest honnête} Contrats à terme, non le comptant~; clôtures synchrones~; la saisonnalité séparée du retour à la moyenne.
\sfield{Sources} Girma et Paulson (1999)~; EIA spot~: ratio de Sharpe de 0.48 (0.35 sans avril 2020)~; ce chapitre~: 1.05 sur un crack construit.
\end{strategyfile}

\begin{strategyfile}{Spark spread avec une vue sur le heat rate}\label{strat:s2:commodity-spreads:spark}
\sfield{Qui paie, et pourquoi} Les producteurs et les détaillants qui couvrent leurs positions~; la vue du marché sur les centrales qui fixent le prix.
\sfield{Instruments et plateformes} Contrats à terme ou forwards sur l'électricité et le gaz, pour la même période et la même zone.
\sfield{Signal} Le heat rate du trader pour la centrale marginale, comparé à celui qu'implique le marché.
\sfield{Dimensionnement et exécution} Gaz au heat rate par mégawattheure d'électricité.
\sfield{Coûts} La large fourchette de l'électricité~; les décalages de profil et de zone.
\sfield{Comment elle meurt} Les nouvelles capacités, les pannes et la météo.
\sfield{Horizon, capacité, infrastructure} De quelques mois à quelques saisons~; un modèle électrique fondamental.
\sfield{Backtest honnête} Le parc de centrales tel qu'il était à l'époque.
\sfield{Sources} Aucune performance vérifiée.
\end{strategyfile}

\begin{strategyfile}{Trade de crush sur le soja}\label{strat:s2:commodity-spreads:crush}
\sfield{Qui paie, et pourquoi} La couverture de marge des transformateurs~; la capacité de trituration referme les marges élargies.
\sfield{Instruments et plateformes} Contrats à terme sur le soja, le tourteau et l'huile (le board crush).
\sfield{Signal} Les écarts par rapport à un équilibre saisonnier~; retour à la moyenne à court terme.
\sfield{Dimensionnement et exécution} 44 livres de tourteau et 11 d'huile par boisseau, en unités de contrat.
\sfield{Coûts} Trois jambes.
\sfield{Comment elle meurt} Les variations de la demande d'huile (biocarburants) ou de tourteau qui déplacent l'équilibre.
\sfield{Horizon, capacité, infrastructure} De quelques jours à quelques mois.
\sfield{Backtest honnête} La tendance de l'équilibre~; les périodes hors échantillon.
\sfield{Sources} Simon (1999).
\end{strategyfile}

\begin{strategyfile}{Spread géographique Brent-WTI}\label{strat:s2:commodity-spreads:location}
\sfield{Qui paie, et pourquoi} Les expéditeurs et les raffineurs qui paient pour accéder à l'autre marché~; le spread rémunère les pipelines et les pétroliers.
\sfield{Instruments et plateformes} Contrats à terme sur le Brent et le WTI.
\sfield{Signal} Le spread comparé aux coûts de transport et à la capacité.
\sfield{Dimensionnement et exécution} Un contrat contre un.
\sfield{Coûts} Faibles~; décalage de roll entre les échéances des deux contrats.
\sfield{Comment elle meurt} Un goulet d'étranglement qui élargit le spread pendant un an~; les corners au point de livraison.
\sfield{Horizon, capacité, infrastructure} De quelques semaines à quelques années~; données de pipelines et d'exportation.
\sfield{Backtest honnête} Prix synchrones~: des évaluations au comptant à des clôtures différentes créent un faux retour à la moyenne.
\sfield{Sources} EIA (2012, 2017)~; EIA spot~: ratio de Sharpe de 0.85, 0.32 avec un jour de retard.
\end{strategyfile}

\begin{strategyfile}{Spread de qualité sweet-sour}\label{strat:s2:commodity-spreads:quality}
\sfield{Qui paie, et pourquoi} Les raffineurs qui peuvent ou non traiter le brut soufré~; les raffineries complexes le paient moins cher.
\sfield{Instruments et plateformes} Différentiels de qualité, swaps ou cargaisons physiques.
\sfield{Signal} Le différentiel comparé à la valeur du traitement supplémentaire.
\sfield{Dimensionnement et exécution} Baril contre baril.
\sfield{Coûts} Différentiels peu liquides.
\sfield{Comment elle meurt} Les variations de l'offre d'une qualité (réductions de l'OPEP sur le brut soufré)~; les conversions de raffineries.
\sfield{Horizon, capacité, infrastructure} Quelques mois.
\sfield{Backtest honnête} Prix évalués, non prix négociés.
\sfield{Sources} Aucune performance vérifiée~; ce chapitre~: 0.98 sur un spread construit.
\end{strategyfile}

\begin{strategyfile}{Crack d'essence saisonnier}\label{strat:s2:commodity-spreads:season}
\sfield{Qui paie, et pourquoi} Le passage à l'essence estivale et la saison de conduite~; la maintenance des raffineurs au printemps.
\sfield{Instruments et plateformes} Essence contre contrats à terme sur le brut, contrats d'été.
\sfield{Signal} Le calendrier.
\sfield{Dimensionnement et exécution} Position acheteuse de février à avril.
\sfield{Coûts} Deux jambes, avec roll.
\sfield{Comment elle meurt} Une saisonnalité bien connue, intégrée à la courbe à terme~; un choc de demande.
\sfield{Horizon, capacité, infrastructure} Trois mois par an.
\sfield{Backtest honnête} Contrats à terme sur le mois d'été, dont le prix intègre déjà la saisonnalité, non le comptant.
\sfield{Sources} EIA spot~: hausse dans 14 années sur 20, moyenne de 6.17~\$, $t$ de 2.8.
\end{strategyfile}

\section{Tutoriel~: le prix du pétrole n'est pas tout}\label{tut:s2:commodity-spreads:tut}

\begin{tutorial}
\textbf{Objectif.} Construire des spreads par ratio, simuler un crack avec saisonnalité et décalage et un \omterm{def:s2:commodity-spreads:location}{spread géographique} avec goulet d'étranglement, parier sur leur retour, et comparer avec
l'historique EIA spot. \textbf{État final~:} le tableau et deux figures.
\begin{tutsteps}
\item \textbf{Spreads par ratio}.
\omcode{code/firm/commspread/firm_commspread.py}{31}{41}{Le crack 3-2-1, le spark spread et le board crush.}\label{lst:s2:commodity-spreads:build}
\item \textbf{Le pari sur le retour et la saisonnalité}.
\omcode{code/firm/commspread/firm_commspread.py}{107}{141}{Parier sur l'écart par rapport à la moyenne glissante, et la règle saisonnière.}\label{lst:s2:commodity-spreads:fade}
\item \textbf{Exécuter} \texttt{table()}, \texttt{bottleneck()}, \texttt{bands()}, \texttt{fig\_commspread.py} et, une fois,
\texttt{s2\_fetch\_spreads.py}.
\end{tutsteps}
\textbf{Pour aller plus loin.} Ajouter des heures de clôture asynchrones au \omterm{def:s2:commodity-spreads:location}{spread géographique} synthétique et mesurer le faux retour à la moyenne~; estimer la saisonnalité
à partir des seules années passées~; arrêter le pari sur le \omterm{def:s2:commodity-spreads:location}{spread géographique} lorsque la capacité de transport est saturée.
\end{tutorial}

\section{Construction~: spreads de matières premières}\label{bld:s2:commodity-spreads:commspread}

\begin{build}
\bfield{Rôle} Construction de spreads par ratio~; un marché synthétique de spreads de brut, de crack, géographiques et de qualité~; règles de pari et règles saisonnières.
\bfield{Interface} \texttt{crack\_321}, \texttt{spark}, \texttt{board\_crush}, \texttt{SpreadConfig(\dots)}, \texttt{simulate\_spreads(cfg)},
\texttt{fade(x, cfg, month, delay)}, \texttt{seasonal(x, month, cfg)}.
\bfield{Règles} P\&L par unité de spread à partir de la variation du lendemain~; coûts par unité négociée~; fenêtres glissantes uniquement.
\bfield{Tests d'acceptation} \texttt{code/firm/commspread/tests/}~: les trois spreads à la main~; les positions et le P\&L du pari à la main~; la fenêtre saisonnière~;
la forme du goulet d'étranglement et la réponse négative du crack au brut.
\bfield{Pour aller plus loin} Courbes à terme pour chaque jambe~; un modèle de raffinerie avec des baisses de run.
\end{build}

\begin{omsources}
\item P.~B.~Girma et A.~S.~Paulson, ``Risk arbitrage opportunities in petroleum futures spreads'', \emph{Journal of Futures Markets} 19(8), 1999.
\item D.~P.~Simon, ``The soybean crush spread: empirical evidence and trading strategies'', \emph{Journal of Futures Markets} 19(3), 1999.
\item US Energy Information Administration, ``2011 brief: energy commodity price trends varied widely during 2011'', \emph{Today in Energy}, 2012.
\item US Energy Information Administration, ``Transportation constraints and export costs widen the Brent-WTI crude oil price spread'',
\emph{Today in Energy}, 15 novembre 2017.
\item US Energy Information Administration, daily spot prices of WTI, Brent, New York Harbor gasoline and diesel, via FRED, consulté le 25 septembre
2026.
\end{omsources}

\section{Exercices}

\begin{exercise}[$\star$]\label{exo:s2:commodity-spreads:1}
Le sans-plomb vaut 2.50~\$ le gallon, le gazole 3.00~\$ et le brut 80~\$ le baril. Calculer le 3-2-1 crack.
\end{exercise}

\begin{exercise}[$\star$]\label{exo:s2:commodity-spreads:2}
L'électricité vaut 50~\$ le mégawatt-heure et le gaz 3~\$ le million de British thermal units. Calculer le spark spread pour un heat rate de 7, puis de 7.5.
\end{exercise}

\begin{exercise}[$\star$]\label{exo:s2:commodity-spreads:3}
Le soja vaut 10~\$ le bushel, le tourteau 300~\$ la short ton et l'huile 50 cents la livre. Calculer le board crush.
\end{exercise}

\begin{exercise}[$\star\star$]\label{exo:s2:commodity-spreads:4}
Pourquoi les variations quotidiennes du crack spread sont-elles corrélées négativement à celles du brut~?
\end{exercise}

\begin{exercise}[$\star\star$]\label{exo:s2:commodity-spreads:5}
Pourquoi le fade Brent-WTI perd-il la plus grande part de son ratio de Sharpe lorsqu'il est initié un jour plus tard~?
\end{exercise}

\begin{exercise}[$\star\star$]\label{exo:s2:commodity-spreads:6}
Quelle force économique rend transitoires les écarts du crush spread~?
\end{exercise}

\begin{exercise}[$\star\star\star$]\label{exo:s2:commodity-spreads:7}
\emph{Programmation.} Exécuter \texttt{bands()}. Comment le ratio de Sharpe du fade sur crack et le temps passé sur le marché évoluent-ils avec des bandes d'entrée de 1.0, 1.5 et 2.0
écarts types, et que donnerait, sur le backtest, le choix a posteriori de la meilleure d'entre elles~?
\end{exercise}

\begin{exercise}[$\star\star\star$]\label{exo:s2:commodity-spreads:8}
\emph{Trouver la faille.} «~Notre fade géographique a rapporté 35~\$ le baril en quatre ans~; le spread revient toujours au coût du transport, alors nous l'avons doublé.~»
\end{exercise}

\section{Problème~: Pas le prix du pétrole}

\begin{problem}[{Problème du week-end --- spreads sur matières premières}]\label{pb:s2:commodity-spreads:1}
Les spreads synthétiques du chapitre et le relevé de l'EIA.

\medskip\noindent\textbf{Partie I --- Spreads.}
\begin{enumerate}
\item Définir un \omterm{def:s2:commodity-spreads:processing}{trade de marge de transformation}, un \omterm{def:s2:commodity-spreads:location}{spread géographique} et un \omterm{def:s2:commodity-spreads:quality}{spread de qualité}.
\item Qu'est-ce qui borne physiquement chaque spread~?
\item Que Girma et Paulson ont-ils établi~?
\item Que Simon a-t-il établi pour le crush~?
\end{enumerate}

\medskip\noindent\textbf{Partie II --- Le relevé.}
\begin{enumerate}[resume]
\item Donner la moyenne, la demi-vie et la corrélation avec le brut du 3-2-1 crack réel.
\item Que s'est-il passé pour le Brent-WTI en 2011 et en avril 2020~?
\item Pourquoi les variations quotidiennes du Brent-WTI reviennent-elles à la moyenne, et qu'est-ce que cela change pour un backtest~?
\item Qu'a fait le crack de l'essence réel de février à avril~?
\end{enumerate}

\medskip\noindent\textbf{Partie III --- Le marché synthétique.}
\begin{enumerate}[resume]
\item Décrire le crack synthétique.
\item Donner le tableau des règles.
\item Pourquoi supprimer la saisonnalité divise-t-elle par deux le fade sur crack~?
\item Qu'a fait le goulot d'étranglement au fade géographique~?
\end{enumerate}

\medskip\noindent\textbf{Partie IV --- Le verdict.}
\begin{enumerate}[resume]
\item Énoncer le \emph{résultat nommé}~: le ratio de Sharpe du trade de crack spread et sa corrélation avec le brut.
\item Pourquoi le ratio de Sharpe de 1.69 de la combinaison est-il trompeur~?
\item Quels sont les trois risques cachés dans un livre de spreads~?
\item Comment backtester honnêtement le fade sur crack~?
\item Quand un fade géographique doit-il être arrêté~?
\item Quelle fiche de stratégie a besoin d'un modèle fondamental~?
\item Quel est le lien entre un livre de spreads et la macro systématique du chapitre~16~?
\item En une phrase~: sur quoi parie un trader de spreads~?
\end{enumerate}
\end{problem}

\section{Questions d'entretien}

\begin{interviewq}[$\star$ \iqroles{trader}]\label{iq:s2:commodity-spreads:1}
Qu'est-ce qu'un crack spread 3-2-1, et qui le couvre~?
\end{interviewq}

\begin{interviewq}[$\star\star$ \iqroles{chercheur}]\label{iq:s2:commodity-spreads:2}
Un spread présente un fort retour à la moyenne dans les données quotidiennes. Comment vérifieriez-vous qu'il ne s'agit pas d'un artefact~?
\end{interviewq}

\begin{interviewq}[$\star\star$ \iqroles{trader}]\label{iq:s2:commodity-spreads:3}
Le Brent-WTI s'est élargi de 10~\$ en un mois. Le fadez-vous~?
\end{interviewq}

\begin{interviewq}[$\star\star$ \iqroles{risque}]\label{iq:s2:commodity-spreads:4}
Comment fixeriez-vous des limites sur un livre de spreads sur matières premières~?
\end{interviewq}

\begin{interviewq}[$\star\star$ \iqroles{développeur}]\label{iq:s2:commodity-spreads:5}
Comment stockeriez-vous et feriez-vous le roll des jambes d'un spread à trois jambes dans un backtester~?
\end{interviewq}

\begin{interviewq}[$\star\star\star$ \iqroles{chercheur}]\label{iq:s2:commodity-spreads:6}
Deux prix suivent la même marche aléatoire mais sont relevés à des moments différents de la journée, séparés d'une fraction $f$ de journée. Montrer que les variations quotidiennes
de leur différence ont une autocorrélation d'ordre un de $-1/2$ quel que soit $f$, et dire ce qui se passe lorsque le spread réel bouge aussi.
\end{interviewq}
